\section{Conclusions}
\label{sec:conclusion}

This report has developed and evaluated a complexity-preserving joint OMP--MP receiver for AFDM under superimposed pilots, augmented by a 1\,041-parameter learned top-$\Kadd$ support-refinement head and evaluated on a doubly-dispersive simulator with first-order MSML modelling. Three conclusions follow directly from the quantitative evidence in Chapter~\ref{sec:results} and the mechanistic interpretation in Chapter~\ref{sec:discussion}, mirroring the three contributions stated in \S\ref{ssec:contrib}.

\begin{itemize}[leftmargin=1.3em,itemsep=4pt,topsep=2pt]

\item \textbf{A complexity-preserving joint OMP--MP receiver for AFDM-SP under fractional-Doppler regimes.}
The receiver runs OMP only once, locks the recovered support, and iterates only the cheap downstream blocks: a damped data-aided least-squares refit of the path gains and sparse MP-BP detection on the chirp-domain factor graph. The single OMP pass already covers $\approx 95\%$ of the Dirichlet-expanded oracle support, with the residual $\approx 5\%$ stable across SNR and therefore structurally rather than noise-driven, so re-running OMP across outer iterations is a $5\times$ constant-factor overhead that adds no support information (Fig.~\ref{fig:iterconv}, Table~\ref{tab:supportrecall}). The receiver thereby recovers two to three orders of magnitude in BER over the one-shot OMP+MP baseline at $\snr = 25$~dB while remaining in the sparse $\mathcal{O}(\Ppath \Nsub^2)$ arithmetic class, two orders of magnitude cheaper than dense LMMSE channel estimation at $\Nsub = 1024$ (Table~\ref{tab:complexity}).

\item \textbf{A 1\,041-parameter learned support-refinement head delivers a consistent $\approx 2$~dB paired BER separation at $\approx 8\%$ MAC overhead.}
A small deep-unfolding head, deployed deterministically as a fixed-budget top-$\Kadd$ ranker between the single OMP pass and the iterative LS/MP-BP loop, scores off-support candidates and appends the highest-ranked atoms to the active support. Pooled across three independent evaluation-seed blocks, the resulting receiver delivers paired BER reductions of $-37.7\%$, $-40.2\%$, and $-37.6\%$ relative to the classical joint at $\snr = 20$, $25$, and $30$~dB respectively, all statistically significant under continuity-corrected paired McNemar (worst-case $p < 10^{-6}$, Table~\ref{tab:ber4b}). The receiver remains in the same $\mathcal{O}(\Ppath \Nsub^2)$ arithmetic class as the classical baseline, with $\approx 8\%$ analytical complex-MAC overhead at $\Nsub = 1024$ (Table~\ref{tab:complexity}; measured wall-clock parity at $\Nsub = 1024$ and $+3.9\%$ at $\Nsub = 2048$). At $\snr \leq 10$~dB the gain is essentially zero ($|\Delta| \leq 6.1\%$), which is the intended behaviour of a support-refinement module in the noise-dominated regime where support errors carry no exploitable structural signature.

\item \textbf{A mechanism analysis grounds the learned gain in measurable support-selection statistics.}
A top-$\Ppath$ physical-support oracle, formed by replacing the OMP step with the simulator's ground-truth main-lobe atoms, is a valid ceiling on the BER reduction available from any support-refinement strategy operating on the proposed dictionary; it reaches the censored-zero regime on $5.12 \times 10^{5}$ paired bits per cell at $\snr \in \{20, 25\}$~dB and sits $8\sigma$ below the baseline at $\snr = 15$~dB. Against this ceiling, the learned head closes $61\%$, $81\%$, and $98\%$ of the baseline-to-oracle gap at $\snr = 15$, $20$, and $25$~dB respectively (\S\ref{ssec:probes}, Table~\ref{tab:oraclehead}). Combined with the support-recall analysis above, this attributes the learned head's BER contribution specifically to the recovery of Dirichlet side-lobe atoms that single-pass OMP discards under fractional-Doppler leakage, rather than to any generic noise-averaging effect.

\end{itemize}

\subsection*{Future work}

Four directions would sharpen the contribution and extend the design statement made here from a simulator-based observation toward an operational underwater acoustic modem claim.

The most immediate is a multi-seed retraining study, with three to five independent training seeds per SNR cell, which would convert the present single-checkpoint observation into a learning-variance statement and place error bars on the pooled $\approx 2$~dB headline. A direct re-implementation of the SP-AFDM channel estimator of Zheng \emph{et al.}~\cite{qing2024spafdm} on the present $(\mathrm{channel}, \mathrm{data}, \mathrm{noise})$ seed blocks would convert the qualitative structural comparison of Table~\ref{tab:priorart} into a quantitative prior-art comparison and resolve whether the OMP-once design loses, matches, or surpasses an SP-AFDM receiver that re-runs the sparse step.

Extending the simulator from the first-order per-path MSML approximation (per-path constant Doppler shift plus delay rescaling) to full broadband per-path waveform resampling and per-frame Doppler jitter, following Cao \emph{et al.}~\cite{zhou2024msmlafdm}, would bring the physical-layer model closer to real UWA propagation and test whether the mechanism-level conclusions identified here transfer to a more realistic propagation regime. Finally, an engineering-deployability study measuring runtime, memory-hierarchy effects, and energy consumption on a reference embedded target such as a DSP or FPGA implementation would extend the complexity-preserving claim from per-frame arithmetic into deployment cost, completing the path from receiver design to operational specification.
